\section{Метод решения}
Площадь треугольника с вершинами $A$, $B$, $C$ в пространстве считается через векторное произведение: $S = \frac{1}{2}|\vec{u} \times \vec{v}|$, где $\vec{u} = B - A$, $\vec{v} = C - A$. При переборе корень не берётся, сравнивается $|\vec{u} \times \vec{v}|^2$: функция $\sqrt{x}/2$ монотонная, поэтому максимум будет у той же тройки, а \texttt{sqrt} вызывается только один раз для ответа.

Перебираются все тройки индексов $i < j < k$, всего их $C_n^3 = \frac{n(n-1)(n-2)}{6}$, сложность $O(n^3)$. Условие $i < j < k$ нужно, чтобы не проверять один треугольник несколько раз и не брать тройки с повторяющейся точкой.

Для распараллеливания работа делится на <<строки>>: одна строка --- это фиксированное $i$ и все $j, k$ после него. Строки независимы, но разные по объёму: в строке $i$ примерно $(n-i)^2/2$ троек, то есть первые строки намного тяжелее последних. Если поделить строки на равные куски подряд, первый поток получит больше половины всей работы. Поэтому строки раздаются динамически: есть общий счётчик \texttt{next\_row}, и каждый свободный поток под мьютексом берёт из него следующую строку. Так быстрые потоки берут больше строк, и все заканчивают примерно одновременно.

Каждый поток ищет максимум в своих строках в локальной переменной без блокировок. Когда строки кончаются, поток один раз под вторым мьютексом сравнивает свой результат с общим \texttt{best} и при необходимости обновляет его. Мьютексы захватываются только на выдаче строк ($n$ раз) и на слиянии результата ($t$ раз), поэтому потоки почти не мешают друг другу. Один поток никогда не держит два мьютекса сразу, поэтому взаимной блокировки быть не может. При равной площади побеждает тройка с меньшим $i$, чтобы ответ не зависел от количества потоков.

Количество потоков задаётся ключом \texttt{-t}: главный поток создаёт ровно $t$ рабочих потоков через \texttt{pthread\_create}, и других потоков во время работы не создаётся, поэтому одновременно работает не больше $t$ потоков. Главный поток сам не считает, а ждёт рабочие потоки через \texttt{pthread\_join}. Поэтому всего в процессе $t + 1$ поток.

Время измеряется через \texttt{clock\_gettime(CLOCK\_MONOTONIC)} только для поиска (создание потоков, перебор, join), генерация точек в замер не входит.

\section{Описание программы}
Проект собирается через CMake (по умолчанию в режиме Release) и разделён на файлы:

\begin{itemize}
\item \texttt{include/points.h}, \texttt{src/points.c} --- структура \texttt{Point} с полями \texttt{x}, \texttt{y}, \texttt{z} и функции:
  \begin{itemize}
  \item \texttt{generate\_points(n, seed)} --- генерирует $n$ случайных точек с координатами в $[-1000, 1000]$;
  \item \texttt{read\_points(filename, \&n)} --- читает точки из файла (в первой строке количество, дальше по три числа в строке);
  \item \texttt{triangle\_area(a, b, c)} --- площадь треугольника.
  \end{itemize}
\item \texttt{include/search.h}, \texttt{src/search.c} --- многопоточный поиск:
  \begin{itemize}
  \item структура \texttt{Triangle} --- индексы $i, j, k$ и значение $|\vec{u} \times \vec{v}|^2$;
  \item структура \texttt{ThreadArg} --- номер потока и сколько строк он посчитал;
  \item \texttt{worker} --- функция потока: берёт строки из счётчика, перебирает тройки, сливает результат;
  \item \texttt{lock}, \texttt{unlock} --- захват и освобождение мьютекса с проверкой кода возврата;
  \item \texttt{find\_max\_triangle(points, n, threads, \&result)} --- создаёт мьютексы и потоки, ждёт их и освобождает ресурсы.
  \end{itemize}
\item \texttt{src/main.c} --- разбор ключей через \texttt{getopt} (\texttt{-t} число потоков, \texttt{-n} количество случайных точек, \texttt{-s} seed, \texttt{-f} файл с точками), проверка аргументов через \texttt{strtol}, замер времени и вывод результата.
\item \texttt{show\_threads.sh} --- запускает программу в фоне и показывает её потоки через \texttt{ps -L}, \texttt{/proc/PID/status} и \texttt{/proc/PID/task}.
\item \texttt{bench.sh} --- замеры времени для разного количества точек и потоков, считает ускорение и эффективность.
\end{itemize}

Используемые системные вызовы и функции POSIX:
\begin{itemize}
\item \texttt{pthread\_create()} (\texttt{clone3}) --- создание потока;
\item \texttt{pthread\_join()} (\texttt{futex}) --- ожидание завершения потока;
\item \texttt{pthread\_mutex\_init()}, \texttt{pthread\_mutex\_destroy()} --- создание и удаление мьютекса;
\item \texttt{pthread\_mutex\_lock()}, \texttt{pthread\_mutex\_unlock()} --- захват и освобождение мьютекса (при ожидании внутри вызывается \texttt{futex});
\item \texttt{gettid()}, \texttt{getpid()} --- получение TID потока и PID процесса для вывода;
\item \texttt{clock\_gettime()} --- замер времени;
\item \texttt{malloc()}, \texttt{free()} (\texttt{mmap}, \texttt{brk}) --- выделение памяти;
\item \texttt{fopen()}, \texttt{fscanf()}, \texttt{fclose()} (\texttt{openat}, \texttt{read}, \texttt{close}) --- чтение файла с точками;
\item \texttt{write()} (через \texttt{printf}) --- вывод.
\end{itemize}

Функции pthread возвращают код ошибки, а не $-1$, и не выставляют \texttt{errno}, поэтому ошибки выводятся через \texttt{strerror(err)}. Если \texttt{pthread\_create} не смог создать поток, join делается только для уже созданных потоков, после чего освобождаются память и мьютексы.
